\section{Related Work}
\label{sec:related}

The study connects offline-to-online RL, transfer across transition dynamics,
and bus holding. We distinguish the algorithmic lineage from the mechanism
that is actually confirmed in the final experiment.

\subsection{Offline and offline-to-online reinforcement learning}

Soft actor-critic (SAC) provides the entropy-regularized, off-policy backbone
used by the online learners in this study~\citep{haarnoja2018soft}. Its later
formulation adds automatic entropy-temperature adjustment
~\citep{haarnoja2018applications}.
Conservative Q-learning (CQL) depresses values for actions outside offline
support~\citep{kumar2020cql}; IQL instead fits an
expectile state value and extracts an advantage-weighted policy without
evaluating unseen actions in its fitted target~\citep{kostrikov2022iql}. We use
IQL to initialize the frozen H2O+ actors and SAC-style updates for their online
continuation.

Offline-to-online fine-tuning can destroy a useful offline policy when early
online data shift the state--action distribution and bootstrap errors
propagate. Balanced replay with a pessimistic critic ensemble is one
stabilization route: it prioritizes online and near-on-policy offline samples
~\citep{lee2022balanced}, rather than simply fixing a 50/50 replay split.
Calibrated Q-learning (Cal-QL) prevents excessive pessimism by requiring
policy-averaged conservative estimates to stay above a reference policy's
value~\citep{nakamoto2023calql}. Policy expansion (PEX) retains an
offline policy and adaptively composes it with a newly learned policy
~\citep{zhang2023pex}. Reinforcement learning with prior data (RLPD) mixes
prior and newly collected replay while learning from random initialization
~\citep{ball2023efficient}, whereas warm-start reinforcement learning (WSRL) begins
from an offline policy, collects a frozen-policy warmup buffer, and then trains
from online replay~\citep{zhou2025wsrl}. These methods motivate the stabilization
controls in the development lineage, but the formal comparison in this paper
does not claim a universal ranking among their source algorithms.

Reward-guided conservative Q-learning (RG-CQL) has combined offline training
and online fine-tuning for coordinating ride-pooling with public transit
~\citep{hu2025rgcql}. We therefore make no broad novelty claim for
offline-to-online RL in transportation. Our focus is an asynchronous fixed-route
holding problem with a target-simulator dynamics gap and an explicit completion
boundary.

\subsection{Transfer across dynamics}

Off-dynamics RL corrects source-domain rewards using a conditional
target-to-source dynamics log ratio estimated by domain classifiers
~\citep{eysenbach2021off}. H2O combines a fixed target-domain
dataset with interaction in an imperfect simulator and uses dynamics-aware
correction when simulator transitions are poorly supported by target dynamics
~\citep{niu2022h2o}. H2O+ broadens that construction across offline and online
base learners~\citep{niu2025h2oplus}.

Our adaptation includes physical event duration in a joint/marginal classifier,
but the final training profile uses unit simulator weights. Thus the classifier
is a domain-shift diagnostic, not the confirmed mechanism. Its empirical
validation is reported in \Cref{sec:exp:validity,app:ratio}; the formal comparison
isolates the action reserve applied to a fixed learned actor.

\subsection{Bus holding and service completion}

The bus-bunching literature spans demand, supply, and control mechanisms
~\citep{rezazada2024review}. \citet{daganzo2009headway} studies headway-based
holding and its stability, whereas \citet{xuan2011dynamic} uses arrival
deviations from a virtual schedule to improve schedule reliability. These are
distinct control objectives. \citet{delgado2012holding} compares holding with
boarding limits under passenger-time objectives. Real-time optimization
incorporates changing demand and running times
~\citep{sanchezmartinez2016realtime}, while predictive control combines
headway forecasts with dynamic holding~\citep{andres2017predictive}.

RL formulations include high-frequency single-service control
~\citep{alesiani2018reinforcement}, multi-agent dynamic holding
~\citep{wang2020real}, asynchronous multi-agent coordination
~\citep{wang2021asynchronous}, and shared-corridor multi-line control
~\citep{wang2023multiobjective}. Asynchronous bus control is therefore not
itself a new contribution. Our study instead combines immutable
target replay, online interaction in a separate reduced-order simulator,
adaptive development of a fixed action transform, and final confirmation on
fresh target-simulator cells. The common passenger ledger and the paired final
reserve comparison distinguish this evaluation from a generic claim that a
new RL learner outperforms existing bus-holding algorithms.

The completion-safety reserve is intentionally simple. It is not a new generic
safe-RL algorithm and does not certify arbitrary operational constraints.
Instead, it encodes one observed transit requirement: when a lightly loaded
tail bus has no follower, a leader exists, and demand remains, suppressing hold
can improve accrued time while leaving passengers unfinished. The paired parent
comparison changes only the reserve, allowing this concrete mechanism to be
tested without reselecting a checkpoint or retraining the actor.
